
\subsection{算子 $\mathfrak p$ 的构造概要}
\label{subsec:3-8-3}

构造 $\mathfrak p$ 的主要想法，是利用空间
$\mathcal M_{k,1}^{\mathrm{main}}(L;\beta)$。这里
$\beta\in\pi_2(M,L)$，$k=0,1,2,\ldots$；下面给出该空间的定义。
令 $\widetilde{\mathcal M}_{k,1}^{\mathrm{main}}(L;\beta)$ 为所有
\[
 (w,\vec z,z^+)=\bigl((\Sigma,\vec z,z^+),w\bigr)
\]
组成的集合，其中
\[
 (\Sigma,\vec z,z^+)\in\mathcal M_{k,1}^{\mathrm{main}},
 \qquad
 \vec z=(z_1,\ldots,z_k),\qquad
 z_i\in\partial\Sigma,
\]
$z^+$ 是一个内部标记点，且
$w\colon\Sigma\to M$ 是满足
\[
 w(\partial\Sigma)\subset L,\qquad [w]=\beta
\]
的伪全纯映射。（$\Sigma$ 是一个圆盘 $D^2$，并且可能带有球面气泡。）
对该集合取 $\mathrm{PSL}(2;\mathbb R)$-商，再加入稳定映射使商空间紧化，
所得空间记作
$\mathcal M_{k,1}^{\mathrm{main}}(L;\beta)$；见定义 2.1.24。
$\mathcal M_{0,1}^{\mathrm{main}}(L;\beta)$ 的紧化需要特别处理，本小节
稍后将作说明。

\begin{proposition}{3.8.21}\label{prop:3-8-21}
若 $k\geq1$，则
$\mathcal M_{k,1}^{\mathrm{main}}(L;\beta)$ 带有带角的仓西结构。
\end{proposition}

注意，$\mathbb Z_k$ 通过循环置换标记点作用在
$\mathcal M_{k,1}^{\mathrm{main}}(L;\beta)$ 上；具体地，
\[
 \operatorname{cyc}
 \bigl(w,(z_1,\ldots,z_k),z^+\bigr)
 =
 \bigl(w,(z_k,z_1,\ldots,z_{k-1}),z^+\bigr).
\]

\begin{proposition}{3.8.22}\label{prop:3-8-22}
$\operatorname{cyc}$ 作用保持仓西结构。它在仓西结构的意义下保持定向，
当且仅当 $k$ 为奇数。\liftIndex{I113}
\end{proposition}

这两个命题的证明与横截性问题密切相关，故推迟到第 7.4 节。扰动仓西
映射来构造虚基本链，同时保持循环对称性，是一件并不平凡的事。

令
\[
\begin{aligned}
 ev=(ev_\partial,ev_0)
 &=(ev_1,\ldots,ev_k,ev_0)\\
 &\colon
 \mathcal M_{k,1}^{\mathrm{main}}(L;\beta)
 \longrightarrow L^k\times M
\end{aligned}
\]
为评价映射，其中
\[
 ev\bigl(w,(z_1,\ldots,z_k),z^+\bigr)
 =
 \bigl(w(z_1),\ldots,w(z_k),w(z^+)\bigr).
\]
设 $P_i=(|P_i|,f_i)$ 是代表 $C(L)$ 某个生成元的奇异链。置
\[
\begin{aligned}
 &\mathcal M_{k,1}^{\mathrm{main}}
 (L;\beta;P_1,\ldots,P_k)\\
 &\quad=
 \mathcal M_{k,1}^{\mathrm{main}}(L;\beta)_{ev_\partial}
 \mathbin{\times_{f_1\times\cdots\times f_k}}
 (P_1\times\cdots\times P_k).
\end{aligned}
\]
这个纤维积带有仓西结构；我们选取适当的多值截面
$\mathfrak s$，从而定义它的虚基本链。构造细节将在第 7.4 节完成。

\begin{definition}{3.8.23}\label{def:3-8-23}
置
\[
 \mathfrak p_{k,\beta}(P_1,\ldots,P_k)
 =
 ev_{0*}\!\left(
 \mathcal M_{k,1}^{\mathrm{main}}
 (L;\beta;P_1,\ldots,P_k)
 \right),
\]
并定义
\begin{equation}
 \mathfrak p_k
 =
 \sum_\beta
 T^{\omega[\beta]}e^{\mu(\beta)/2}\mathfrak p_{k,\beta}.
\tag{3.8.24}\label{eq:3-8-24}
\end{equation}
\end{definition}

\eqref{eq:3-8-24} 右端的收敛性照常由格罗莫夫紧致性得到。

\begin{lemma}{3.8.25}\label{lem:3-8-25}
可以选取多值截面 $\mathfrak s$，使
$\mathfrak p_{k,\beta}$ 满足下列循环性质：
\[
 \mathfrak p_{k,\beta}(P_1,\ldots,P_k)
 =
 (-1)^{
  \deg'P_k\left(\sum_{i=1}^{k-1}\deg'P_i\right)}
 \mathfrak p_{k,\beta}(P_k,P_1,\ldots,P_{k-1}).
\]
\end{lemma}

本书不证明算子 $\mathfrak m_k$ 具有类似的循环对称性，见注 3.8.8。
引理 3.8.25 将在第 7.4 节证明，相关定向问题则在第 8.10.2 小节分析。

现在证明 \eqref{eq:3-8-10-1}--\eqref{eq:3-8-10-3}。
\eqref{eq:3-8-10-1} 很容易。事实上，$\beta_0=0$ 时的模空间
$\mathcal M_{1,1}^{\mathrm{main}}(L;\beta_0)$ 已经横截，并且可与 $L$
本身等同。因此无须扰动它，且
\[
 \mathfrak p_{1,\beta_0}(P)=P,
\]
这就推出 \eqref{eq:3-8-10-1}。

下面证明 \eqref{eq:3-8-10-2}。为了得到
$\delta_M\circ\mathfrak p_k$ 一项，考察
$\mathcal M_{k,1}^{\mathrm{main}}(L;\beta;P_1,\ldots,P_k)$ 的边界。
若
\[
 \left(
  (w^{(m)},\vec z^{\,(m)},z^{+(m)}),
  (p_1^{(m)},\ldots,p_k^{(m)})
 \right)
\]
是一列发散序列，则发生下列情形之一。
\translatornote{原书在这列纤维积点的末尾又重复写了一个
$z^{+(m)}$；该内部标记点已经属于前面的模空间元素，并不是纤维积的另一
分量，译文据对象类型删去重复项。}

\smallskip
\noindent\textup{(3.8.26.1)}
某个 $p_i^{(m)}$ 趋向 $P_i$ 的边界。

\noindent\textup{(3.8.26.2)}
$w^{(m)}\colon\Sigma\to M$ 趋向一个映射
\[
 w^{(\infty)}\colon\Sigma_1\cup\Sigma_2\longrightarrow M,
\]
其中两个亏格零有边黎曼曲面沿边界粘合。

\smallskip
当然也可能出现球面气泡，但它不贡献边界，因为球面气泡出现在余维二；
研究稳定映射时，它总是出现在实余维二。
情形 (3.8.26.1) 对应
\[
 \mathfrak p_{k,\beta}
 (P_1,\ldots,\partial P_i,\ldots,P_k),
\]
其中 $\partial$ 是 $C(L)$ 上的经典边界算子，并满足
$\mathfrak m_{1,0}=(-1)^n\partial$；见定义 3.5.6。

下面研究情形 (3.8.26.2)。令
\[
 \bigl((p_1^{(\infty)},\ldots,p_k^{(\infty)}),z^{+(\infty)}\bigr)
\]
为各标记点的极限。可给 $\Sigma_1,\Sigma_2$ 编号，使
$z^{+(\infty)}\in\Sigma_2$。分成以下三种情形。

\smallskip
\noindent\textup{(3.8.26.2.1)}
存在 $a\leq b$，使
$p_a^{(\infty)},\ldots,p_b^{(\infty)}$ 位于 $\partial\Sigma_1$，
而其余 $p_i^{(\infty)}$ 位于 $\partial\Sigma_2$。

\noindent\textup{(3.8.26.2.2)}
存在 $a\leq b$，使
\[
 p_{b+1}^{(\infty)},\ldots,p_k^{(\infty)},
 p_1^{(\infty)},\ldots,p_{a-1}^{(\infty)}
\]
位于 $\partial\Sigma_1$，而其余 $p_i^{(\infty)}$ 位于
$\partial\Sigma_2$。

\noindent\textup{(3.8.26.2.3)}
$\partial\Sigma_1$ 上没有任何 $p_i^{(\infty)}$。

\smallskip
见图 3.8.1--3.8.3。情形 (3.8.26.2.1) 对应
\[
 \mathfrak p
 \bigl(P_1,\ldots,\mathfrak m(P_a,\ldots,P_b),\ldots,P_k\bigr);
\]
情形 (3.8.26.2.2) 对应
\[
 \mathfrak p
 \bigl(
 P_a,\ldots,P_b,
 \mathfrak m(P_{b+1},\ldots,P_k,P_1,\ldots,P_{a-1})
 \bigr);
\]
情形 (3.8.26.2.3) 对应
\[
 \mathfrak p
 \bigl(P_1,\ldots,P_i,\mathfrak m_0(1),P_{i+1},\ldots,P_k\bigr).
\]

\begin{figure}[htbp]
  \centering
  \setlength{\unitlength}{0.85mm}
  \begin{picture}(126,66)(0,-7)
    \put(39,29){\circle{42}}
    \put(81,29){\circle{42}}
    \put(60,29){\circle*{1.7}}
    \qbezier[18](28,54)(7,48)(8,27)
    \qbezier[16](8,27)(9,10)(29,5)
    \qbezier[22](68,54)(94,63)(108,42)
    \qbezier[18](108,42)(116,27)(106,13)
    \qbezier[16](106,13)(92,1)(69,5)
    \put(84,31){\circle*{1.7}}
    \put(87,30){$z^{+(\infty)}$}
    \put(28,49){\line(1,-1){2.5}}
    \put(17,40){\line(1,-1){2.5}}
    \put(31,9){\line(1,1){2.5}}
    \put(48,13){\line(-1,1){2.5}}
    \put(26,51){$p_a^{(\infty)}$}
    \put(2,42){$p_{a+1}^{(\infty)}$}
    \put(20,1){$p_{b-1}^{(\infty)}$}
    \put(45,5){$p_b^{(\infty)}$}
    \put(68,48){\line(1,-1){2.5}}
    \put(101,42){\line(-1,-1){2.5}}
    \put(102,16){\line(-1,1){2.5}}
    \put(68,10){\line(1,1){2.5}}
    \put(61,51){$p_{a-1}^{(\infty)}$}
    \put(106,43){$p_1^{(\infty)}$}
    \put(106,10){$p_k^{(\infty)}$}
    \put(57,1){$p_{b+1}^{(\infty)}$}
    \put(30,-5){$\partial\Sigma_1$}
    \put(74,-5){$\partial\Sigma_2$}
  \end{picture}
  \par\smallskip
  \renewcommand{\thefigure}{3.8.1}
  \refstepcounter{figure}
  \textbf{\figurename~\thefigure}\label{fig:3-8-1}
\end{figure}

\begin{figure}[htbp]
  \centering
  \setlength{\unitlength}{0.85mm}
  \begin{picture}(126,66)(0,-7)
    \put(39,29){\circle{42}}
    \put(81,29){\circle{42}}
    \put(60,29){\circle*{1.7}}
    \qbezier[18](28,54)(7,48)(8,27)
    \qbezier[16](8,27)(9,10)(29,5)
    \qbezier[22](68,54)(94,63)(108,42)
    \qbezier[18](108,42)(116,27)(106,13)
    \qbezier[16](106,13)(92,1)(69,5)
    \put(84,31){\circle*{1.7}}
    \put(87,30){$z^{+(\infty)}$}
    \put(28,49){\line(1,-1){2.5}}
    \put(17,31){\line(1,0){2.5}}
    \put(48,13){\line(-1,1){2.5}}
    \put(49,45){\line(-1,-1){2.5}}
    \put(19,51){$p_k^{(\infty)}$}
    \put(1,29){$p_1^{(\infty)}$}
    \put(38,4){$p_{a-1}^{(\infty)}$}
    \put(43,50){$p_{b+1}^{(\infty)}$}
    \put(68,48){\line(1,-1){2.5}}
    \put(68,10){\line(1,1){2.5}}
    \put(65,51){$p_b^{(\infty)}$}
    \put(63,2){$p_a^{(\infty)}$}
    \put(30,-5){$\partial\Sigma_1$}
    \put(74,-5){$\partial\Sigma_2$}
  \end{picture}
  \par\smallskip
  \renewcommand{\thefigure}{3.8.2}
  \refstepcounter{figure}
  \textbf{\figurename~\thefigure}\label{fig:3-8-2}
\end{figure}

\begin{figure}[htbp]
  \centering
  \setlength{\unitlength}{0.92mm}
  \begin{picture}(126,63)
    \put(38,31){\circle{30}}
    \put(76,31){\circle{46}}
    \put(53,31){\circle*{1.7}}
    \qbezier[22](64,57)(95,67)(111,43)
    \qbezier[18](111,43)(119,27)(108,11)
    \put(76,33){\circle*{1.7}}
    \put(79,31){$z^{+(\infty)}$}
    \put(62,51){\line(1,-1){2.5}}
    \put(62,12){\line(1,1){2.5}}
    \put(98,43){\line(-1,-1){2.5}}
    \put(99,18){\line(-1,1){2.5}}
    \put(52,53){$p_i^{(\infty)}$}
    \put(47,4){$p_{i+1}^{(\infty)}$}
    \put(102,44){$p_1^{(\infty)}$}
    \put(103,11){$p_k^{(\infty)}$}
    \put(25,7){$\partial\Sigma_1$}
    \put(68,2){$\partial\Sigma_2$}
  \end{picture}
  \par\smallskip
  \renewcommand{\thefigure}{3.8.3}
  \refstepcounter{figure}
  \textbf{\figurename~\thefigure}\label{fig:3-8-3}
\end{figure}

计入引理 3.8.25 的循环对称性后，用与定理 3.5.11 的证明相同的论证，
便得到 \eqref{eq:3-8-10-2}。

现在证明 \eqref{eq:3-8-10-3}。为此需要指出下列现象。考虑正合序列
\[
 H_2(M)\xrightarrow{i_*}H_2(M,L)
 \xrightarrow{\partial}H_1(L).
\]
当 $\partial\beta=0\in H_1(L)$ 时，为了紧化
$\mathcal M_{0,1}(L;\beta)$，必须纳入不稳定映射。

假设 $\partial\beta=0$，则对某个
$\widetilde\beta\in H_2(M)$ 有
$\beta=i_*(\widetilde\beta)$；这样的 $\widetilde\beta$ 可能不止一个。
考虑满足
\[
 (w,(z_0^+,z_1^+))
 \in\mathcal M_{0,2}(M;\widetilde\beta),
 \qquad w(z_0^+)\in L
\]
的元素。不要把
$\mathcal M_{0,2}(M;\widetilde\beta)$ 与
$\mathcal M_{0,2}(L;\beta)$ 的记号混淆：前者下标 $0,2$ 分别表示定义域
（闭黎曼曲面）的亏格为 $0$ 和内部标记点数为 $2$；后者下标 $0,2$
则分别表示边界标记点数为 $0$ 和内部标记点数为 $2$。

考虑映射
\[
 \operatorname{const}_{w(z_0^+)}\colon D^2\longrightarrow M,
 \qquad
 \operatorname{const}_{w(z_0^+)}(x)\equiv w(z_0^+).
\]
在 $0\in D^2$ 与 $z_0^+\in S^2$ 处把 $D^2$ 和 $S^2$ 粘合，得到
$\Sigma$；再考虑
\[
 \bigl((\Sigma,z_1^+),
 w\mathbin{\#}\operatorname{const}_{w(z_0^+)}\bigr).
\]
这不是通常意义下的稳定映射，因为
\[
 \operatorname{Aut}(D^2,0)\cong U(1)
\]
并非有限群。然而，这种对象会作为
$\mathcal M_{0,1}(L;\beta)$ 中元素的极限出现；见第 7.4 节稍后给出的
命题 3.8.27 的证明。这类对象与
\[
 \mathcal M_{0,2}(M;\widetilde\beta)\mathbin{\times_M}L
\]
一一对应。综上可得：

\begin{proposition}{3.8.27}\label{prop:3-8-27}
考虑
\[
 \mathcal M_{0,1}(L;\beta)
 \ \cup\
 \bigcup_{\widetilde\beta\,;\,i_*(\widetilde\beta)=\beta}
 \left(
 \mathcal M_{0,2}(M;\widetilde\beta)\mathbin{\times_M}L
 \right).
\]
可以在这个空间上定义一个拓扑，使其成为紧豪斯多夫空间。它带有定向的
带边界仓西结构。其中
$\mathcal M_{0,2}(M;\widetilde\beta)\mathbin{\times_M}L$
位于其边界，而该纤维积上的仓西结构与
\[
 \mathcal M_{0,1}(L;\beta)
 \ \cup\
 \bigcup_{\widetilde\beta}
 \left(
 \mathcal M_{0,2}(M;\widetilde\beta)\mathbin{\times_M}L
 \right)
\]
上的仓西结构相容。
\end{proposition}

本命题将在第 7.4.1 小节证明。关键在于，此处的球面气泡（对应
$\mathcal M_{0,2}(M;\widetilde\beta)\mathbin{\times_M}L$）出现在余维一。
这是因为带该球面气泡的常值圆盘具有一维自同构群。

以下为简化记号，以
\[
 \mathcal M_{0,1}(L;\beta)
 \cup
 \left(
 \mathcal M_{0,2}(M;(i_*)^{-1}(\beta))
 \mathbin{\times_M}L
 \right)
\]
表示
\[
 \mathcal M_{0,1}(L;\beta)
 \cup
 \bigcup_{\widetilde\beta}
 \left(
 \mathcal M_{0,2}(M;\widetilde\beta)\mathbin{\times_M}L
 \right).
\]
利用这个模空间，按定义 3.8.23 的同样方式定义
$\mathfrak p_{0,\beta}$。

现在可以证明 \eqref{eq:3-8-10-3}。上述模空间的边界是
\[
 \mathcal M_{0,2}(M;(i_*)^{-1}(\beta))
 \mathbin{\times_M}L
\]
与
\[
 \sum_{\beta_1+\beta_2=\beta}
 \mathcal M_1(L;\beta_1)
 \mathbin{\times_L}
 \mathcal M_{1,1}(L;\beta_2)
\]
之并。因此
\[
\begin{aligned}
 &\delta_M\bigl(\mathfrak p_{0,\beta}(1)\bigr)
 +\sum_{\beta_1+\beta_2=\beta}
 \mathfrak p_{1,\beta_2}
 \bigl(\mathfrak m_{0,\beta_1}(1)\bigr)\\
 &\qquad
 +\sum_{\widetilde\beta\,;\,i_*(\widetilde\beta)=\beta}
 ev_{0*}\left(
 \mathcal M_{0,2}(M;\widetilde\beta)
 \mathbin{\times_M}L
 \right)=0.
\end{aligned}
\]
这便推出 \eqref{eq:3-8-10-3}。式
\eqref{eq:3-8-10-4}--\eqref{eq:3-8-10-6} 的证明，则是第 7.2 节与
第 7.3 节论证的结合。
\translatornote{原书此处只列 (3.8.10.4)、(3.8.10.5)；但定理 3.8.9
前后的说明以及第 7.4 节完成证明处都明确把 (3.8.10.6) 与二者并列，
译文据此补足连续范围。}

\subsection{算子 $\mathfrak q$}
\label{subsec:3-8-4}

我们通过把最后一个（第 $0$ 个）边界标记点移入内部，构造了算子
$\mathfrak p$。现在把第 $1,\ldots,k$ 个边界标记点移入内部，从而构造
算子 $\mathfrak q$。\liftIndex{I038}\liftIndexRangeStart{I357}
正如下个小节将讨论的，这一构造使定义弗洛尔上同调所需的条件得以放宽；
例如，它可用于证明导言中的定理 I。还要指出，
\textbf{[Ono96, Sei02, OzSa04]} 也利用环境辛流形中的闭链截取伪全纯
圆盘模空间。

为了精确陈述本小节的主要技术结果定理 3.8.32，先准备一些记号。以下均取
$R=\mathbb Q$。像第 3.5 节那样，取 $L$ 上光滑奇异单形的可数集
$\mathcal X_1(L)$，令
$C_1(L;\Lambda_{0,\mathrm{nov}})$ 为由 $\mathcal X_1(L)$ 生成的上链
复形。第 3.5 节的滤过 $A_\infty$ 结构 $\mathfrak m$ 正是构造在这个
复形上。

令 $S(M)$ 为 $M$ 的光滑奇异链复形，并把它视为上链复形。选取
$M$ 上光滑奇异单形的一个可数集 $\mathcal X(M)$，令 $C(M)$ 为
$S(M)$ 中由 $\mathcal X(M)$ 生成的 $\mathbb Q$-子模，并置
\[
 C(M;\Lambda_{0,\mathrm{nov}})
 =
 C(M)\widehat\otimes\Lambda_{0,\mathrm{nov}}.
\]

\begin{statement}{假设}{3.8.28}\label{assump:3-8-28}
\noindent\textup{(3.8.29.1)}
$C(M)$ 是 $S(M)$ 的上链子复形。

\noindent\textup{(3.8.29.2)}
自然包含 $C(M)\to S(M)$ 诱导上同调同构
\[
 H(C(M),\delta_M)\cong H(M).
\]

\noindent\textup{(3.8.29.3)}
每个 $Q\in\mathcal X(M)$ 都与 $L$ 横截，且 $L\cap Q$ 具有光滑三角
剖分；该剖分中的每个单形都属于 $\mathcal X_1(L)$。

\noindent\textup{(3.8.29.4)}
此外，$\partial_iQ\cap L$ 的三角剖分与 $Q\cap L$ 的三角剖分之限制
一致，其中 $\partial_iQ$ 是 $Q$ 的第 $i$ 个面。
\end{statement}

固定 $L\cap Q$ 的三角剖分，使 (3.8.29.3) 成立，并定义
\[
 i^\pitchfork\colon C(M)\longrightarrow C(L),
 \qquad i^\pitchfork(Q)=L\cap Q.
\]
于是 $L\cap Q$ 确定 $C(L)$ 的一个元素。由 (3.8.29.4)，
$i^\pitchfork$ 是上链同态；它显然诱导由包含映射给出的自然同态
$H^*(M)\to H^*(L)$。

把 $C(M;\Lambda_{0,\mathrm{nov}})$ 中元素的次数移位 $2$，得到复形
$C(M;\Lambda_{0,\mathrm{nov}})[2]$，并考虑相伴的 bar 复形
\[
 BC(M;\Lambda_{0,\mathrm{nov}})[2].
\]
如下定义对称群 $\mathfrak S_k$ 在
$B_kC(M;\Lambda_{0,\mathrm{nov}})[2]$ 上的作用：
\[
 \sigma(y_1\otimes\cdots\otimes y_k)
 =
 (-1)^{\epsilon(\sigma)}
 y_{\sigma(1)}\otimes\cdots\otimes y_{\sigma(k)}
 \qquad(\sigma\in\mathfrak S_k),
\]
其中
\[
 \epsilon(\sigma)
 =
 \sum_{\substack{i<j\\\sigma(i)>\sigma(j)}}
 \deg y_i\,\deg y_j.
\]
这里次数移位了 $2$，所以无须使用移位次数。

\begin{definition}{3.8.30}\label{def:3-8-30}
$E_kC(M;\Lambda_{0,\mathrm{nov}})[2]$ 是
$BC(M;\Lambda_{0,\mathrm{nov}})[2]$ 中由 $\mathfrak S_k$-作用不变元
组成的子空间。置
\[
 \widehat EC(M;\Lambda_{0,\mathrm{nov}})[2]
 =
 \widehat{\bigoplus}_k
 E_kC(M;\Lambda_{0,\mathrm{nov}})[2].
\]
\liftIndex{I096}\liftIndex{I139}
\end{definition}

$BC(M;\Lambda_{0,\mathrm{nov}})[2]$ 上的余积 $\Delta$ 在
$\widehat EC(M;\Lambda_{0,\mathrm{nov}})[2]$ 上诱导余积：
\begin{equation}
\begin{aligned}
 \Delta\colon
 \widehat EC(M;\Lambda_{0,\mathrm{nov}})[2]
 \longrightarrow{}&
 \widehat EC(M;\Lambda_{0,\mathrm{nov}})[2]\\
 &{}\widehat\otimes
 \widehat EC(M;\Lambda_{0,\mathrm{nov}})[2].
\end{aligned}
\tag{3.8.31}\label{eq:3-8-31}
\end{equation}
也就是说，把 $\Delta$ 限制在
$\widehat EC(M;\Lambda_{0,\mathrm{nov}})[2]$ 上，其像包含于
\eqref{eq:3-8-31} 的右端。$C(M;\Lambda_{0,\mathrm{nov}})$ 上的上边界
算子 $\delta_M$ 在 $\widehat EC(M;\Lambda_{0,\mathrm{nov}})[2]$ 上
诱导余导子 $\widehat\delta_M$：
\[
\begin{aligned}
 &\widehat\delta_M(y_1\otimes\cdots\otimes y_k)\\
 &\qquad=
 \sum_{i=1}^k
 (-1)^{\deg y_1+\cdots+\deg y_{i-1}}
 y_1\otimes\cdots\otimes
 \delta_My_i\otimes\cdots\otimes y_k.
\end{aligned}
\]
见第 A3 节。

令
$y\in EC(M)[2]$、
$\mathbf x\in BC(L)[1]$，并置
\[
 \Delta^{k-1}y
 =
 \sum_{c_1}
 y_{c_1}^{(k;1)}\otimes\cdots\otimes y_{c_1}^{(k;k)},
 \qquad
 \Delta^{k-1}\mathbf x
 =
 \sum_{c_2}
 \mathbf x_{c_2}^{(k;1)}\otimes\cdots
 \otimes\mathbf x_{c_2}^{(k;k)}.
\]
本节始终采用这一记号。\liftIndex{I038}

本小节的主要结果如下。

\begin{theorem}{3.8.32}\label{thm:3-8-32}
存在包含 $\mathcal X_1(L)$ 的光滑奇异单形可数集
$\mathcal X_{1+}(L)$，并且由 $\mathcal X_{1+}(L)$ 生成的
$C_{1+}(L;\Lambda_{0,\mathrm{nov}})$ 上存在一个滤过 $A_\infty$ 代数
结构，使
$(C_{1+}(L;\Lambda_{0,\mathrm{nov}}),\mathfrak m)$ 具有下列性质。

\begin{enumerate}[label=\textup{(\arabic*)},leftmargin=2.8em]
  \item $C_{1+}(L;\Lambda_{0,\mathrm{nov}})$ 上的滤过 $A_\infty$ 代数
  结构延拓
  $C_1(L;\Lambda_{0,\mathrm{nov}})$ 上的滤过 $A_\infty$ 代数结构。

  \item 存在次数为 $+1$ 的一列算子
  \[
  \begin{aligned}
   \mathfrak q_{\ell,k}\colon{}&
   E_\ell C(M;\Lambda_{0,\mathrm{nov}})[2]
   \otimes
   B_kC_{1+}(L;\Lambda_{0,\mathrm{nov}})[1]\\
   &\longrightarrow C_{1+}(L;\Lambda_{0,\mathrm{nov}})[1],
  \end{aligned}
  \]
  其中 $k,\ell=0,1,2,\ldots$，并满足
  \begin{equation}
  \begin{aligned}
   &\mathfrak q_{\ell,k}
   \bigl(\widehat\delta_M(y)\otimes\mathbf x\bigr)\\
   &\quad+
   \sum_{c_1,c_2}
   (-1)^{
    \deg y_{c_1}^{(2;2)}
      \deg'\mathbf x_{c_2}^{(3;1)}
    +\deg'\mathbf x_{c_2}^{(3;1)}
    +\deg y_{c_1}^{(2;1)}}\\
   &\qquad\cdot
   \mathfrak q\left(
    y_{c_1}^{(2;1)}
    \otimes
    \left(
     \mathbf x_{c_2}^{(3;1)}
     \otimes
     \mathfrak q\bigl(
      y_{c_1}^{(2;2)}
      \otimes\mathbf x_{c_2}^{(3;2)}
     \bigr)
     \otimes\mathbf x_{c_2}^{(3;3)}
    \right)
   \right)=0.
  \end{aligned}
  \tag{3.8.33}\label{eq:3-8-33}
  \end{equation}
  若 $\ell\ne0$，则
  \begin{equation}
   \mathfrak q_{\ell,k}
   \equiv
   \begin{cases}
    i^\pitchfork\pmod{\Lambda^+_{0,\mathrm{nov}}},
      &\ell=1,\ k=0,\\
    0\pmod{\Lambda^+_{0,\mathrm{nov}}},
      &\text{其他情形}.
   \end{cases}
  \tag{3.8.34.1}\label{eq:3-8-34-1}
  \end{equation}
  当 $\ell=0$ 时，在
  $B_kC_{1+}(L;\Lambda_{0,\mathrm{nov}})[1]$ 上有
  \begin{equation}
   \mathfrak q_{0,k}=\mathfrak m_k.
  \tag{3.8.34.2}\label{eq:3-8-34-2}
  \end{equation}

  \item 还可把 $\mathfrak q_{\ell,k}$ 延拓为
  \[
  \begin{aligned}
   \mathfrak q^+_{\ell,k}\colon{}&
   E_\ell C(M;\Lambda_{0,\mathrm{nov}})[2]
   \otimes
   B_kC_{1+}(L;\Lambda_{0,\mathrm{nov}})^+[1]\\
   &\longrightarrow
   C_{1+}(L;\Lambda_{0,\mathrm{nov}})^+[1].
  \end{aligned}
  \]
  在 \eqref{eq:3-8-33} 中以 $\mathfrak q^+$ 代替 $\mathfrak q$ 后，
  该式仍成立，并且
  \begin{align}
   \mathfrak q^+_{\ell,k}
   \bigl(y\otimes
    (\mathbf x_1\otimes\mathbf e^+\otimes\mathbf x_2)\bigr)
   &=0
   \qquad(\ell\geq1),
   \tag{3.8.35.1}\label{eq:3-8-35-1}\\
   \mathfrak q^+_{0,k}
   &=\mathfrak m_k^+.
   \tag{3.8.35.2}\label{eq:3-8-35-2}
  \end{align}

  \item 此外，置
  \[
   \widetilde{\mathbf e}=\operatorname{PD}[M]
  \]
  （即基本类的庞加莱对偶），并置
  \[
  \begin{aligned}
   C(M;\Lambda_{0,\mathrm{nov}})^+
   ={}&C(M;\Lambda_{0,\mathrm{nov}})\\
   &{}\oplus
   \Lambda_{0,\mathrm{nov}}[\widetilde{\mathbf f}]
   \oplus
   \Lambda_{0,\mathrm{nov}}[\widetilde{\mathbf e}^+].
  \end{aligned}
  \]
  这里
  $\deg\widetilde{\mathbf f}=-1$、
  $\deg\widetilde{\mathbf e}^+=0$。于是可以进一步把
  $\mathfrak q^+_{\ell,k}$ 延拓为
  \[
  \begin{aligned}
   \mathfrak q^{++}_{\ell,k}\colon{}&
   E_\ell C(M;\Lambda_{0,\mathrm{nov}})^+[2]
   \otimes
   B_kC_{1++}(L;\Lambda_{0,\mathrm{nov}})^+[1]\\
   &\longrightarrow
   C_{1++}(L;\Lambda_{0,\mathrm{nov}})^+[1].
  \end{aligned}
  \]
  在 \eqref{eq:3-8-33} 中以 $\mathfrak q^{++}$ 代替
  $\mathfrak q$，并在 (3.8.35) 各式中以 $\mathfrak q^{++}$ 代替
  $\mathfrak q^+$ 后，各式仍成立。此外，当
  $(\ell,k)\ne(0,0)$ 时，
  \begin{equation}
   \mathfrak q^{++}_{\ell+1,k}
   \bigl(y\otimes\widetilde{\mathbf e}^+,\mathbf x\bigr)=0,
  \tag{3.8.36.1}\label{eq:3-8-36-1}
  \end{equation}
  而且\translatornote{原书 (3.8.36.3) 的右端写作
  $\mathbf e^+-\mathbf e$；但 $\delta_M$ 在本项刚定义的环境复形
  $C(M;\Lambda_{0,\mathrm{nov}})^+$ 上作用，而无波浪号的两个幺元属于
  $L$ 上的复形。译文依该式的型别及紧前定义恢复为两个带波浪号的环境
  幺元。}
  \begin{align}
   \mathfrak q^{++}_{1,0}
   (\widetilde{\mathbf e}^+,1)
   &=\mathbf e^+,
   \tag{3.8.36.2}\label{eq:3-8-36-2}\\
   \delta_M(\widetilde{\mathbf f})
   &=\widetilde{\mathbf e}^+-\widetilde{\mathbf e},
   \tag{3.8.36.3}\label{eq:3-8-36-3}\\
   \mathfrak q^{++}_{1,0}
   (\widetilde{\mathbf f},1)
   &\equiv\mathbf f\pmod{\Lambda^+_{0,\mathrm{nov}}}.
   \tag{3.8.36.4}\label{eq:3-8-36-4}
  \end{align}
\end{enumerate}
\end{theorem}
\liftIndex{I038}\liftIndexRangeEnd{I357}

定理 3.8.32 的证明将在第 3.8.6 小节概述，完整证明见第 7.4 节和
第 7.3.3 小节。

注意，$\mathfrak q^+_{\ell,k}$ 来自定义 3.3.2 中 $\mathfrak h$ 的延拓
\translatornote{原书误写成“定义 3.3.3”；$\mathfrak h_k$ 实际定义于
定义 3.3.2，而 (3.3.3) 是该定义内的公式 $\mathfrak h_1=\mathfrak m$。}
\[
\begin{aligned}
 \mathfrak h_k^{+M}\colon{}&
 EC(M;\Lambda_{0,\mathrm{nov}})[2]
 \otimes
 B_k\bigl(BC_{1+}(L;\Lambda_{0,\mathrm{nov}})[1]\bigr)\\
 &\longrightarrow C_{1+}(L;\Lambda_{0,\mathrm{nov}})[1],
\end{aligned}
\]
其关系为
\begin{equation}
\begin{aligned}
 &\mathfrak q^+\bigl(
 y\otimes(
 \mathbf x_1\otimes\mathbf f\otimes
 \mathbf x_2\otimes\mathbf f\otimes\cdots
 \otimes\mathbf f\otimes\mathbf x_k)
 \bigr)\\
 &\qquad=
 \mathfrak h_k^{+M}
 \bigl(
 y,\mathbf x_1\mathbin{\overline\otimes}\cdots
 \mathbin{\overline\otimes}\mathbf x_k
 \bigr),
\end{aligned}
\tag{3.8.35.3}\label{eq:3-8-35-3}
\end{equation}
但 $\ell=0$、$k=2$ 且
$\mathbf x_1,\mathbf x_2\in\Lambda_{0,\mathrm{nov}}$ 时除外。在这一
例外情形下，
\begin{equation}
\begin{aligned}
 \mathfrak q^+_{0,1}(1\otimes\mathbf f)
 &=\mathfrak m_1^+(\mathbf f)\\
 &=\mathbf e^+-\mathbf e
 +\mathfrak h_2^{+M}
 \bigl(1,1\mathbin{\overline\otimes}1\bigr),
\end{aligned}
\tag{3.8.35.4}\label{eq:3-8-35-4}
\end{equation}
这与 \eqref{eq:3-8-35-2} 相容。这里
$C_{1+}(L;\Lambda_{0,\mathrm{nov}})^+[1]$、$\mathbf e^+$ 与
$\mathbf f$ 均如第 3.3 节。

\eqref{eq:3-8-33} 是一个相当复杂的公式。第 7.4 节将引入若干代数
结构（$L_\infty$ 结构），并用它们重写定理 3.8.32；所得结果就是导言的
定理 Y 和定理 7.4.102。相关代数讨论还可参见
\textbf{[KaSt05]}。
